% TODO checklist:
% - [x] Inspected src/workers/jobs_worker.py (483 lines - complete worker implementation)
% - [x] Inspected db/redis_cache/jobs_cache.py (queue operations)
% - [x] Inspected src/services/jobs_service.py (service layer)
% - [x] Inspected Q&A.md Q39 for worker implementation details
% - [x] Analyzed signal handling and graceful shutdown patterns

\section{Worker Architecture}
\label{sec:worker-arch}

The worker architecture implements a separate process model for job execution, decoupling the FastAPI web application from long-running computations. This design enables horizontal scaling of worker capacity independently from API throughput, prevents job execution from blocking HTTP request handling, and provides isolation for resource-intensive operations.

This section details the worker process design, queue polling mechanisms, heartbeat monitoring, and graceful shutdown procedures.

\subsection{Process Separation Model}

The Cineca Agentic Platform employs a process-per-role architecture where the web application and job workers run as separate OS processes:

\begin{itemize}
    \item \textbf{Web Application Process}: Handles HTTP requests via FastAPI/Uvicorn, manages job creation, and serves SSE event streams. Started via \texttt{uvicorn src.app:app}.
    
    \item \textbf{Worker Process}: Polls Redis queues for pending jobs, executes job logic, and updates PostgreSQL with results. Started via \texttt{python -m src.workers.jobs\_worker}.
\end{itemize}

This separation provides several benefits:

\begin{enumerate}
    \item \textbf{Resource Isolation}: CPU-intensive job execution cannot starve HTTP request processing.
    
    \item \textbf{Independent Scaling}: Worker replicas can be added without affecting API capacity.
    
    \item \textbf{Fault Isolation}: A worker crash does not bring down the API server.
    
    \item \textbf{Deployment Flexibility}: Workers can run on different node types (e.g., GPU-enabled nodes for ML workloads).
\end{enumerate}

\subsection{Worker Implementation Details}
\label{subsec:worker-impl-details}

The \texttt{JobsWorker} class in \texttt{src/workers/jobs\_worker.py} implements the complete worker lifecycle:

\begin{lstlisting}[language=Python, caption={JobsWorker class initialization.}]
class JobsWorker:
    """Background worker for processing jobs with PostgreSQL persistence."""

    def __init__(
        self,
        poll_interval: float = 1.0,
        heartbeat_interval: float = 5.0,
        max_iterations: int | None = None,
    ):
        self.poll_interval = poll_interval
        self.heartbeat_interval = heartbeat_interval
        self.max_iterations = max_iterations  # For testing
        self.running = False
        self.current_job_id: str | None = None

        # Setup signal handlers for graceful shutdown
        signal.signal(signal.SIGTERM, self._signal_handler)
        signal.signal(signal.SIGINT, self._signal_handler)

    def _signal_handler(self, signum, frame):
        """Handle shutdown signals gracefully."""
        logger.info(f"Received signal {signum}, initiating graceful shutdown...")
        self.running = False
\end{lstlisting}

\subsection{Queue Polling Mechanism}

Workers poll Redis queues using the blocking \texttt{BRPOP} operation, which efficiently waits for new jobs without busy-waiting:

\begin{lstlisting}[language=Python, caption={Main worker loop with queue polling.}]
async def start(self):
    """Main worker loop."""
    logger.info("Jobs worker starting...")
    self.running = True
    iteration = 0

    while self.running:
        # Check max iterations (for testing)
        if self.max_iterations and iteration >= self.max_iterations:
            logger.info(f"Reached max iterations ({self.max_iterations})")
            break

        iteration += 1

        try:
            processed = await self._process_next_job()
            if not processed:
                # No jobs available, sleep before next poll
                await asyncio.sleep(self.poll_interval)

        except Exception as e:
            logger.error(f"Error in worker loop: {e}", exc_info=True)
            await asyncio.sleep(self.poll_interval)

    logger.info("Jobs worker stopped")
\end{lstlisting}

The queue polling iterates through allowed job types in round-robin fashion:

\begin{lstlisting}[language=Python, caption={Round-robin queue polling across job types.}]
async def _process_next_job(self) -> bool:
    """Process one job from the queue."""
    allowed_types_str = getattr(settings, "ALLOWED_JOB_TYPES", "demo,test,long-running")
    allowed_types = [t.strip() for t in allowed_types_str.split(",")]

    # Try to pop from each queue
    for job_type in allowed_types:
        job_id = await asyncio.to_thread(
            jobs_cache.queue_pop_job, job_type, timeout=0
        )

        if job_id:
            logger.info(f"Popped job {job_id} from queue '{job_type}'")
            self.current_job_id = job_id

            try:
                db_generator = get_db()
                db = next(db_generator)
                try:
                    await self._execute_job(job_id, db)
                finally:
                    with contextlib.suppress(StopIteration):
                        next(db_generator)
                    self.current_job_id = None

            except Exception as e:
                logger.error(f"Failed to execute job {job_id}: {e}")

            return True

    return False  # All queues empty
\end{lstlisting}

\subsection{Job Execution Flow}

When a job is dequeued, the worker executes a well-defined sequence of operations:

\begin{enumerate}
    \item \textbf{Load Job}: Retrieve the full job document from PostgreSQL.
    
    \item \textbf{Check Cancellation}: Verify the Redis cancel flag before starting execution.
    
    \item \textbf{Transition to RUNNING}: Atomically update status and log the transition event.
    
    \item \textbf{Execute with Heartbeat}: Run job logic while maintaining periodic heartbeat updates.
    
    \item \textbf{Check Cancellation}: Verify cancel flag after execution completes.
    
    \item \textbf{Transition to Terminal}: Update status to FINISHED, FAILED, or CANCELLED based on outcome.
\end{enumerate}

\begin{lstlisting}[language=Python, caption={Complete job execution flow.}]
async def _execute_job(self, job_id: str, db: Session):
    """Execute a single job with full lifecycle management."""
    job_uuid = UUID(job_id)
    jobs_service = JobsService(db)

    try:
        # Load job from PostgreSQL
        job = jobs_service.repo.get_job(job_uuid)
        if not job:
            logger.error(f"Job {job_id} not found in database")
            return

        # Check if already cancelled
        if await asyncio.to_thread(jobs_cache.check_cancel_flag, job_id):
            logger.info(f"Job {job_id} was cancelled before execution")
            await self._mark_cancelled(job_uuid, jobs_service)
            return

        # Transition to RUNNING
        job = jobs_service.transition_status(
            job_uuid, from_status="queued", to_status="running"
        )

        # Log event for SSE streaming
        jobs_service.append_event(
            job_uuid,
            event_type="status",
            event_data={
                "to": "running", "from": "queued",
                "timestamp": datetime.utcnow().isoformat()
            },
        )

        # Execute job with heartbeat monitoring
        result = await self._run_job_with_heartbeat(
            job_id, job.type, job.payload_json or {}, jobs_service
        )

        # Check if cancelled during execution
        if await asyncio.to_thread(jobs_cache.check_cancel_flag, job_id):
            await self._mark_cancelled(job_uuid, jobs_service)
            return

        # Transition to FINISHED
        job = jobs_service.repo.update_job_result(job_uuid, result)
        job = jobs_service.transition_status(
            job_uuid, from_status="running", to_status="finished"
        )

        jobs_service.append_event(
            job_uuid, event_type="status",
            event_data={"to": "finished", "from": "running", ...}
        )

    except Exception as e:
        logger.error(f"Job {job_id} failed: {e}", exc_info=True)
        # Transition to FAILED
        jobs_service.repo.update_job_error(job_uuid, str(e))
        jobs_service.transition_status(
            job_uuid, from_status="running", to_status="failed"
        )
\end{lstlisting}

\subsection{Heartbeat Mechanism}

The heartbeat mechanism serves two purposes: indicating worker liveness and enabling detection of stale jobs from crashed workers. A concurrent asyncio task updates the job's \texttt{updated\_at} timestamp periodically:

\begin{lstlisting}[language=Python, caption={Heartbeat loop for worker liveness indication.}]
async def _run_job_with_heartbeat(
    self, job_id: str, job_type: str,
    payload: dict, jobs_service: JobsService
) -> dict:
    """Execute job logic with periodic heartbeat updates."""
    job_uuid = UUID(job_id)

    # Start heartbeat task
    heartbeat_task = asyncio.create_task(
        self._heartbeat_loop(job_uuid, jobs_service)
    )

    try:
        result = await self._execute_job_type(job_id, job_type, payload)
        return result
    finally:
        # Stop heartbeat
        heartbeat_task.cancel()
        with contextlib.suppress(asyncio.CancelledError):
            await heartbeat_task

async def _heartbeat_loop(self, job_uuid: UUID, jobs_service: JobsService):
    """Periodically update job timestamp."""
    while True:
        await asyncio.sleep(self.heartbeat_interval)
        try:
            jobs_service.repo.touch_job(job_uuid)
            logger.debug(f"Heartbeat for job {job_uuid}")
        except Exception as e:
            logger.warning(f"Heartbeat failed for job {job_uuid}: {e}")
\end{lstlisting}

The heartbeat interval is configurable via \texttt{JOB\_WORKER\_HEARTBEAT\_INTERVAL} (default: 5 seconds). Jobs with stale heartbeats (no update for multiple intervals) can be detected and recovered by monitoring systems.

\subsection{Job Type Dispatching}

The worker dispatches execution to type-specific handlers based on the \texttt{job.type} field:

\begin{lstlisting}[language=Python, caption={Job type dispatch logic.}]
async def _execute_job_type(
    self, job_id: str, job_type: str, payload: dict
) -> dict:
    """Execute job logic based on job type."""
    if job_type == "demo":
        return await self._execute_demo_job(job_id, payload)
    elif job_type == "test":
        return await self._execute_test_job(job_id, payload)
    elif job_type == "long-running":
        return await self._execute_long_running_job(job_id, payload)
    elif job_type == "agent.run":
        return await self._execute_agent_run_job(job_id, payload)
    else:
        raise ValueError(f"Unknown job type: {job_type}")
\end{lstlisting}

\subsubsection{Demo Job Implementation}

The demo job simulates work with configurable duration while checking for cancellation:

\begin{lstlisting}[language=Python, caption={Demo job with cancellation checking.}]
async def _execute_demo_job(self, job_id: str, payload: dict) -> dict:
    """Execute demo job (sleep simulation)."""
    duration_ms = payload.get("duration_ms", 1000)
    duration_sec = duration_ms / 1000.0

    # Simulate work with cancellation checks every 0.5s
    start = time.time()
    sleep_chunks = int(duration_sec / 0.5) + 1

    for _ in range(sleep_chunks):
        if await asyncio.to_thread(jobs_cache.check_cancel_flag, job_id):
            raise asyncio.CancelledError("Job cancelled")
        await asyncio.sleep(min(0.5, duration_sec))

    elapsed = time.time() - start

    return {
        "status": "completed",
        "requested_duration_ms": duration_ms,
        "actual_duration_ms": int(elapsed * 1000),
        "message": "Demo job completed successfully",
    }
\end{lstlisting}

\subsubsection{Long-Running Job Implementation}

Long-running jobs execute multiple steps with progress updates:

\begin{lstlisting}[language=Python, caption={Long-running job with step-based progress.}]
async def _execute_long_running_job(self, job_id: str, payload: dict) -> dict:
    """Execute long-running job with progress updates."""
    steps = payload.get("steps", 10)
    step_duration = 3.0  # 3 seconds per step

    for step in range(1, steps + 1):
        if await asyncio.to_thread(jobs_cache.check_cancel_flag, job_id):
            raise asyncio.CancelledError("Job cancelled")

        logger.info(f"Job {job_id}: step {step}/{steps}")
        await asyncio.sleep(step_duration)

    return {
        "status": "completed",
        "steps_completed": steps,
        "total_duration_ms": int(steps * step_duration * 1000),
        "message": f"Completed {steps} steps",
    }
\end{lstlisting}

\subsubsection{Agent.Run Job Implementation (Workflow B)}
\label{subsubsec:agent-run-job}

The \texttt{agent.run} job type represents the most complex handler, providing full orchestrator integration within the job system. This handler enables Workflow B execution (see Section~\ref{subsubsec:workflow-b}), bridging the asynchronous job infrastructure with the synchronous orchestration engine.

\paragraph{Implementation Overview}

The \texttt{\_execute\_agent\_run\_job()} method implements a comprehensive execution flow:

\begin{enumerate}
    \item \textbf{AgentRun Record Management}: Load or create the linked \texttt{AgentRun} PostgreSQL record.
    \item \textbf{Orchestrator Initialization}: Instantiate the orchestrator via \texttt{Orchestrator.from\_env()} with full configuration.
    \item \textbf{Parameter Construction}: Build orchestration parameters including temperature, max steps, and principal context.
    \item \textbf{Progress Events}: Emit SSE events at key stages (\texttt{agent\_run\_started}, \texttt{orchestrator\_init}, \texttt{orchestration\_start}).
    \item \textbf{Orchestration Execution}: Execute \texttt{orchestrator.run()} with timeout and cancellation checks.
    \item \textbf{Security Processing}: Apply PII scrubbing and output guard validation.
    \item \textbf{Metrics Emission}: Record agent-specific Prometheus metrics.
    \item \textbf{State Synchronization}: Update both \texttt{Job} and \texttt{AgentRun} records to terminal states.
\end{enumerate}

\paragraph{Payload Schema}

The \texttt{AgentRunJobPayload} Pydantic model validates job payloads:

\begin{lstlisting}[language=Python, caption={AgentRunJobPayload schema for agent.run jobs.}]
class AgentRunJobPayload(BaseModel):
    """Payload schema for agent.run job type."""
    # Required fields
    prompt: str = Field(..., description="The user's prompt/goal")
    user_id: str = Field(..., description="User ID from JWT")
    tenant_id: str = Field(..., description="Tenant identifier")
    
    # Optional fields
    session_id: str | None = None
    run_id: str | None = None
    model: str | None = None
    temperature: float = Field(default=0.2, ge=0.0, le=2.0)
    max_steps: int = Field(default=8, ge=1, le=50)
    metadata: dict[str, Any] | None = None
    principal: dict[str, Any] | None = None
\end{lstlisting}

\paragraph{Progress Events}

The handler emits progress events enabling real-time client updates:

\begin{itemize}
    \item \texttt{agent\_run\_started}: Initial acknowledgment with run ID
    \item \texttt{orchestrator\_init}: Orchestrator ready for execution
    \item \texttt{orchestration\_start}: Beginning orchestration pipeline
    \item \texttt{orchestration\_step}: Per-step progress (from orchestrator callbacks)
    \item \texttt{orchestration\_complete}: Final result with output and metrics
\end{itemize}

\paragraph{Security Integration}

The handler applies two security processing stages to the orchestration output:

\begin{itemize}
    \item \textbf{PII Scrubbing}: Sanitizes personal information from output text and structured data using the \texttt{PIIScrubber} service.
    \item \textbf{Output Guard}: Validates output content against safety policies, filtering or rejecting prohibited content.
\end{itemize}

\paragraph{Bidirectional State Synchronization}

The handler maintains consistency between \texttt{Job} and \texttt{AgentRun} records:

\begin{lstlisting}[language=Python, caption={AgentRun state synchronization.}]
# Update AgentRun with final state
agent_run.status = (
    "succeeded" if result.success else "failed"
)
agent_run.result = result.to_dict()
agent_run.steps_json = [s.to_dict() for s in result.steps]
agent_run.todos_json = result.todos
agent_run.finished_at = datetime.utcnow()
db.commit()

# Emit final progress event
await self._emit_progress_event(
    job_id, "orchestration_complete",
    {"run_id": str(run_id), "success": result.success}
)
\end{lstlisting}

\subsection{Graceful Shutdown}

The worker implements graceful shutdown to ensure clean process termination without leaving jobs in inconsistent states:

\begin{enumerate}
    \item \textbf{Signal Reception}: SIGTERM or SIGINT triggers the signal handler, setting \texttt{self.running = False}.
    
    \item \textbf{Current Job Completion}: The worker completes any in-progress job before exiting the main loop.
    
    \item \textbf{Resource Cleanup}: Database sessions are properly closed via context management.
    
    \item \textbf{Worker Deregistration}: The worker removes itself from the active worker pool (if implemented).
\end{enumerate}

\begin{lstlisting}[language=Python, caption={Signal handler for graceful shutdown.}]
def _signal_handler(self, signum, frame):
    """Handle shutdown signals gracefully."""
    logger.info(f"Received signal {signum}, initiating graceful shutdown...")
    self.running = False
    # The main loop will exit after completing the current job
\end{lstlisting}

The configurable shutdown timeout (via \texttt{JOB\_WORKER\_SHUTDOWN\_TIMEOUT}) allows operators to balance between quick restarts and job completion guarantees.

\subsection{Worker Entry Point}

The worker module provides a standalone entry point for process startup:

\begin{lstlisting}[language=Python, caption={Worker main entry point.}]
async def main():
    """Main entry point for worker process."""
    poll_interval = float(getattr(settings, "JOB_WORKER_POLL_INTERVAL", 1.0))
    heartbeat_interval = float(getattr(settings, "JOB_WORKER_HEARTBEAT_INTERVAL", 5.0))

    # Verify PostgreSQL backend is enabled
    use_postgres = str(getattr(settings, "USE_POSTGRES_JOBS", False)).lower()
    if use_postgres not in ("true", "1", "yes"):
        logger.error("Worker requires USE_POSTGRES_JOBS=true")
        sys.exit(1)

    logger.info("Starting PostgreSQL-backed jobs worker")
    
    worker = JobsWorker(
        poll_interval=poll_interval,
        heartbeat_interval=heartbeat_interval
    )

    try:
        await worker.start()
    except KeyboardInterrupt:
        logger.info("Worker interrupted by user")
    except Exception as e:
        logger.error(f"Worker failed: {e}", exc_info=True)
        sys.exit(1)

if __name__ == "__main__":
    logging.basicConfig(level=logging.INFO)
    asyncio.run(main())
\end{lstlisting}

\subsection{Worker Configuration Parameters}

Table~\ref{tab:worker-config} summarizes the configurable parameters for worker processes.

\begin{table}[ht]
\centering
\caption{Worker configuration parameters.}
\label{tab:worker-config}
\small
\begin{tabular}{lll}
\hline
\textbf{Parameter} & \textbf{Default} & \textbf{Description} \\
\hline
\texttt{JOB\_WORKER\_POLL\_INTERVAL} & 1.0 & Queue polling interval (seconds) \\
\texttt{JOB\_WORKER\_HEARTBEAT\_INTERVAL} & 5.0 & Heartbeat update interval (seconds) \\
\texttt{USE\_POSTGRES\_JOBS} & false & Enable PostgreSQL job persistence \\
\texttt{ALLOWED\_JOB\_TYPES} & demo,test,long-running,agent.run & Permitted job types \\
\texttt{DATABASE\_URL} & - & PostgreSQL connection string \\
\texttt{REDIS\_URL} & - & Redis connection string \\
\hline
\end{tabular}
\end{table}

